\documentclass[12pt,a4paper]{article}
% ===== Packages =====
\usepackage[utf8]{inputenc}
\usepackage[T1]{fontenc}
\usepackage{authblk}
\usepackage[english]{babel}
\usepackage{amsmath, amssymb, amsthm}
\usepackage{mathtools}
\usepackage{algorithm}
\usepackage{algpseudocode}
\usepackage{booktabs}
\usepackage{hyperref}
\usepackage{geometry}
\usepackage{xcolor}
\usepackage{graphicx}
\graphicspath{{../}{./}}
\geometry{margin=2.5cm}

% ===== Macros =====
\newcommand{\Dw}{D_{\mathrm{w}}}
\newcommand{\Ddw}{D}
\newcommand{\Dpv}{D_{\mathrm{PV}}}
\newcommand{\Cinv}{C^{-1}}
\newcommand{\Cpvinv}{C_{\mathrm{PV}}^{-1}}
\newcommand{\Aop}{A}
\newcommand{\PL}{P_{L}}
\newcommand{\Pp}{P_{+}}
\newcommand{\Pm}{P_{-}}
\newcommand{\dg}{^{\dagger}}
\newcommand{\norm}[1]{\left\lVert #1 \right\rVert}

\title{\textbf{A Pauli--Villars Preconditioned, Aggregation-Based\\
Algebraic Multigrid Solver for the M\"obius Domain-Wall Operator}}
\author{W.-L.~Chen}
\affil{\small (working note, Bridge++ \texttt{with\_alpha} branch)}
\date{\today}

\begin{document}
\maketitle

\begin{abstract}
We document the operator and the solution algorithm used in the present
Bridge++ implementation of an algebraic multigrid (AMG) solver for the M\"obius
domain-wall (DW) fermion, together with all parameters of the test problem. The
bare DW operator $\Ddw$ is non-normal and its spectrum straddles the imaginary
axis, so Krylov methods applied directly to $\Ddw$ converge poorly. Following
Kanamori, Chen and Matsufuru [PoS(LATTICE2024)414] we instead solve the
Pauli--Villars (PV) and site-diagonal preconditioned operator
$\Aop=(\Dpv\,\Cpvinv)\dg\,\Cinv\,\Ddw$, whose spectrum lies in a small disk in
the open right half-plane. A left-preconditioning identity makes the solution of
$\Aop$ return $\Ddw^{-1}b$ exactly. The outer iteration is a flexible GMRES
(FGMRES) on $\Aop$, preconditioned by a two-grid V-cycle (a Schwarz
alternating-procedure smoother and an $\Aop$-Galerkin coarse-grid correction
built from chirality-split, block-orthonormalised near-null vectors). We give the
complete operator definitions, the 5D--4D chiral projection of source and sink,
the algorithm, the full set of test parameters, and a \emph{residual-matched}
iteration-count comparison against the production even-odd CGNE solver and against
unpreconditioned GMRES on $\Aop$. We emphasise that the near-null-space setup is
deliberately minimal in this run; the comparison therefore understates what the
method can reach, and is reported as a correctness-and-feasibility baseline.
\end{abstract}

\section{The M\"obius domain-wall operator}
\label{sec:operator}

We work with the five-dimensional M\"obius domain-wall operator $\Ddw$ with a
Wilson kernel. Let $\Dw(M_0)$ be the four-dimensional Wilson--Dirac operator with
negative mass (domain-wall height) $-M_0$, written through the hopping parameter
$\kappa = 1/(8-2M_0)$. $\Dw$ is an \emph{eight-point} nearest-neighbour stencil,
\begin{equation}
(\Dw\psi)(x) = \psi(x) - \kappa\sum_{\mu=1}^{4}\Big[
(1-\gamma_\mu)U_\mu(x)\psi(x+\hat\mu)
+ (1+\gamma_\mu)U_\mu\dg(x-\hat\mu)\psi(x-\hat\mu)\Big],
\end{equation}
i.e.\ two hops ($\pm\hat\mu$) in each of the four spacetime directions per site.
Let $s=0,\dots,N_s-1$ index the fifth dimension of extent $N_s$. With
$P_\pm=\tfrac12(1\pm\gamma_5)$ the M\"obius operator parameterised by $(b,c)$ acts
on a five-dimensional spinor $\psi_s$ as
\begin{equation}
(\Ddw\,\psi)_s
= \big(b\,\Dw + 1\big)\,\psi_s
\;+\; \big(c\,\Dw - 1\big)\!\left[ \Pm \,\psi_{s+1} + \Pp\,\psi_{s-1} \right],
\end{equation}
with the mass-carrying wrap $\psi_{N_s} = -m_q \Pm \psi_0$,
$\psi_{-1} = -m_q \Pp \psi_{N_s-1}$. The fifth-dimension coupling carries no gauge
links, so \textbf{one application of $\Ddw$ costs $N_s$ four-dimensional
eight-point Wilson hops} (plus the cheap tridiagonal $\pm\gamma_5$ shift). The
Pauli--Villars operator is the same operator at unit mass,
$\Dpv \equiv \Ddw\big|_{m_q\to1}$.

The DW operator carries an exactly site-diagonal (in the four-dimensional sense)
factor $C$ in the fifth dimension --- the $N_s\times N_s$ tridiagonal block that
is local in spacetime. Its inverse $\Cinv$ is computed \emph{exactly} by
forward/back substitution (LU of the $N_s$-block, the Bridge++ \texttt{Prec} mode
$L^{-1}\!\circ U^{-1}$); $\Cpvinv$ is the same for the PV operator. These factors
are site-local: they cost \emph{no} four-dimensional hop. No iterative solver and
no double-precision promotion is used for any of them.

The parameters of the test problem are collected in Table~\ref{tab:params}.

\begin{table}[h]
\centering
\begin{tabular}{lll}
\toprule
quantity & symbol & value \\
\midrule
lattice volume        & $L^3\times T$    & $12^4$ \\
gauge action          &                  & quenched Wilson, $\beta=5.7$ \\
configuration         &                  & quenched HMC, plaquette $0.560$ \\
fifth-dimension extent& $N_s$            & $8$ \\
domain-wall height    & $M_0$            & $1.0$ \\
M\"obius coefficients & $(b,c)$          & $(1.5,\,0.5)$ \\
quark mass            & $m_q$            & $0.005$ \\
PV mass               & $m_{\mathrm{PV}}$& $1.0$ \\
alpha parameter       & $\alpha$         & $1.0$ \\
boundary conditions   &                  & $(1,1,1,-1)$ \\
\bottomrule
\end{tabular}
\caption{Operator and gauge parameters of the test problem.}
\label{tab:params}
\end{table}

\subsection{5D--4D chiral projection of source and sink}
\label{sec:projection}
The physical (four-dimensional) propagator is carried on the two chiral boundary
walls. A four-dimensional point source $\eta(x)$ is lifted to the five-dimensional
right-hand side by placing its $\Pp$ and $\Pm$ chiral halves on the two walls and
multiplying by the M\"obius factor $(c\Dw-1)$:
\begin{equation}
b_s \;=\; (c\,\Dw - 1)\,\Big[\,\Pp\,\eta\;\delta_{s,0}
\;+\; \Pm\,\eta\;\delta_{s,N_s-1}\,\Big].
\label{eq:source}
\end{equation}
After solving $\Ddw\,\psi=b$, the four-dimensional propagator is reconstructed by
the chiral projection of the two boundary slices,
\begin{equation}
q(x) \;=\; \Pm\,\psi_{0}(x) \;+\; \Pp\,\psi_{N_s-1}(x)
\;=\; \tfrac12\big[(1-\gamma_5)\psi_{0} + (1+\gamma_5)\psi_{N_s-1}\big].
\label{eq:sink}
\end{equation}
The residual mass is obtained from the same construction applied at the fifth-
dimension \emph{midpoint} cut $s=N_s/2-1,\,N_s/2$,
$q_{\mathrm{mid}} = \Pm\,\psi_{N_s/2} + \Pp\,\psi_{N_s/2-1}$, which defines the
midpoint pseudoscalar density $J_{5q}$. This chiral source/sink construction is
identical across the three solvers compared below, so their pseudoscalar
correlators $C_{PP}(t)$ are directly comparable.

\section{The preconditioned target operator}
\label{sec:target}

The bare operator $\Ddw$ is non-normal: its eigenvalues have both signs of real
part, and unpreconditioned GMRES/FGMRES stagnate. We therefore form the positive
(right-half-plane) target operator of Ref.~[PoS(LATTICE2024)414],
\begin{equation}
\boxed{\;\Aop \;=\; \big(\Dpv\,\Cpvinv\big)\dg\;\Cinv\;\Ddw\;}
\label{eq:A}
\end{equation}
in which all four factors are applied exactly. Empirically (16$^4$, Arnoldi) the
spectrum of $\Aop$ lies in a small disk with \emph{positive real part} and
spectral radius $\approx0.47$, in agreement with the reference. Since $\Cinv$ and
$\Cpvinv$ are site-local, \textbf{one application of $\Aop$ costs two
four-dimensional hopping sweeps} (one through $\Ddw$, one through $\Dpv\dg$), i.e.\
$2N_s = 16$ eight-point Wilson hops on the present lattice.

\paragraph{Left preconditioning.}
Define the left-preconditioner
$\PL = (\Dpv\,\Cpvinv)\dg\,\Cinv$, so that $\Aop = \PL\,\Ddw$. To solve the
physical system $\Ddw\,x = b$ we solve instead
\begin{equation}
\Aop\,x = b', \qquad b' = \PL\,b ,
\label{eq:leftprec}
\end{equation}
because $\Aop x = \PL b \iff \PL(\Ddw x - b)=0$ and $\PL$ is non-singular; the
solution is $x = \Ddw^{-1}b$ exactly. The inner and physical residuals are related
by $b'-\Aop x = \PL(b-\Ddw x)$. \textbf{This relation is essential to a fair
comparison: $\PL$ is ill-conditioned, so a small residual on $\Aop$ does not imply
a small physical residual $\norm{b-\Ddw x}$.} We therefore always drive and report
the \emph{physical} relative residual $\norm{b-\Ddw x}/\norm{b}$ when comparing
solvers, never the $\Aop$-residual alone.

\section{The multigrid algorithm}
\label{sec:algorithm}

The solver is an aggregation-based two-grid algebraic multigrid acting on the
positive operator $\Aop$. The outer Krylov method is \emph{flexible} GMRES, which
is required because the V-cycle preconditioner is not a fixed linear operator.

\subsection{Aggregates, near-null space and the transfer operators}
The fine lattice is partitioned into non-overlapping blocks (aggregates) of size
$3^4$ in spacetime (the full $N_s$ is kept in the fifth dimension), giving a
coarse lattice of $4^4$ sites on the $12^4$ volume. A set of $N_v=8$ near-null
vectors $\{w_k\}$ is generated in the setup phase. Within each aggregate the
vectors are split by chirality $P_\pm$ and \emph{block orthonormalised}; the
chirality split doubles the coarse degrees of freedom to $2N_v=16$ per aggregate.
This defines the prolongation $P$ (piecewise span of the orthonormal block
vectors) and the restriction $R = P\dg$.

\paragraph{Block orthonormalisation.}
The block Gram--Schmidt is performed by a \emph{batched classical} Gram--Schmidt
(CGS): for each target vector all block inner products against the
already-accepted vectors are computed in one batched kernel, subtracted in one
batched \texttt{axpy}, and the block is normalised. A single CGS pass is used. A
second (re-)orthogonalisation pass is \emph{not} applied: the projector basis
carries two chiralities ($\norm{w_k}^2_{\text{block}}=2$, one from $\Pm$ and one
from $\Pp$), and a naive second pass mixes the chirality sectors and destroys the
orthogonality established by the first pass. The single batched CGS pass
reproduces the modified-Gram--Schmidt result to working precision at $O(1)$ kernel
launches per target vector.

\paragraph{Coarse operator and its stencil locality.}
The coarse-grid operator is the Galerkin projection of $\Aop$,
$\Aop_{c} = R\,\Aop\,P = P\dg\,\Aop\,P$, assembled explicitly so the coarse solve
never touches the fine grid. The coarse problem is solved approximately by
restarted GMRES (loose tolerance $10^{-2}$, at most two iterations), which is
adequate inside a V-cycle.

The defining design axis among domain-wall multigrid methods is the
\emph{hopping range (stencil locality) of the coarse operator}, which sets the
cost of building and applying $\Aop_c$. The squared red-black operator used by
HDCG produces a coarse stencil reaching all points of taxicab distance up to
four --- a $321$-point four-dimensional stencil (only reduced to $81$ points at a
minimum $4^4$ block), so re-evaluating the little Dirac operator along an HMC
trajectory costs thousands of fine multiplies. The recent multiple-right-hand-side
domain-wall multigrid likewise adopts a \emph{four-hop} coarse operator. At the
opposite end, a \emph{zero-hop} (purely block-diagonal, no inter-aggregate
coupling) coarse approximation is the cheapest but discards the long-wavelength
coupling that gives multigrid its mass/volume independence. \textbf{Our coarse
operator is a one-hop (nearest-neighbour) coarse Dirac operator}
(\texttt{AFopr\_Domainwall\_coarse}): the Galerkin projection of the
PV-preconditioned $\Aop$ is truncated to a nearest-neighbour stencil, which
preserves sparsity (and hence recursive coarsenability) at a small fraction of
the assembly cost of the four-hop stencils, while retaining the inter-aggregate
coupling absent from the zero-hop choice. A general comparison of these
domain-wall multigrid coarsening strategies is given by Boyle and Yamaguchi
(arXiv:2103.05034).

\subsection{Smoother}
The smoother is a Schwarz alternating procedure (SAP) on the same $3^4$ blocks: a
domain-decomposed block solve of $\Aop$ that damps the high-frequency error the
coarse grid cannot represent. Four smoother sweeps are used, each block solve run
to $10^{-3}$. (A CGNR-on-$\Ddw$ smoother diverges at small $m_q$ with $M_0=1$, the
near-critical regime the PV preconditioning is designed to avoid; the SAP smoother
on $\Aop$ is robust there.)

\subsection{The V-cycle and the outer iteration}
One application of the preconditioner $M^{-1}\approx\Aop^{-1}$ is a two-grid
V-cycle:

\begin{algorithm}[H]
\caption{Two-grid V-cycle $M^{-1}$ applied to a vector $r$ (returns $z$)}
\begin{algorithmic}[1]
\State $z \gets \text{SAP}(\Aop, r)$ \Comment{pre-smoothing, 4 sweeps}
\State $r_c \gets R\,(r - \Aop\,z)$ \Comment{restrict the fine residual}
\State $e_c \gets \Aop_c^{-1} r_c$ \Comment{coarse solve (loose GMRES)}
\State $z \gets z + P\,e_c$ \Comment{prolongate the coarse correction}
\State $z \gets z + \text{SAP}(\Aop,\, r - \Aop\,z)$ \Comment{post-smoothing, 4 sweeps}
\end{algorithmic}
\end{algorithm}

\noindent The complete solver is

\begin{algorithm}[H]
\caption{PV-preconditioned AMG solve of $\Ddw\,x=b$}
\begin{algorithmic}[1]
\State \textbf{Setup:} generate $N_v=8$ near-null vectors; chirality-split and
       block-orthonormalise (CGS); assemble $P,R,\Aop_c$.
\State $b' \gets \PL\,b$ \Comment{left-precondition: now solve $\Aop x = b'$}
\State Solve $\Aop\,x = b'$ by FGMRES (restart length $N_M=30$) using the V-cycle
       $M^{-1}$ as preconditioner, driving the \emph{physical} residual
       $\norm{b-\Ddw x}/\norm{b}$ to the target tolerance.
\State \Return $x = \Ddw^{-1}b$ \Comment{exact by the left-prec identity}
\end{algorithmic}
\end{algorithm}

\noindent All grid-transfer and coarse-grid work is in single precision
(\texttt{float}); the outer FGMRES and the operator $\Aop$ are applied in double.
No multiword/extended precision is used in this study.

\paragraph{Setup parameters and a caveat.}
Table~\ref{tab:mg} lists the multigrid parameters. \textbf{The near-null-space
setup is deliberately minimal here: the number of setup smoothing steps is
$N_{\mathrm{step}}=0$}, so the $N_v=8$ vectors are used essentially as generated,
without the iterative enrichment (applying the smoother to random vectors) that a
tuned aggregation setup performs. The coarse space is therefore \emph{not}
adapted to the near-null modes of $\Aop$, which is expected to leave a sizeable
part of the achievable speed-up on the table. Optimising the setup
($N_{\mathrm{step}}>0$, more vectors $N_v$, possibly more grid levels) is the
natural next step and requires more device memory and throughput than the single
consumer GPU used here.

\begin{table}[h]
\centering
\begin{tabular}{lll}
\toprule
component & parameter & value \\
\midrule
aggregate (block) & spacetime $\times$ fifth dim & $3^4 \times N_s$ \\
coarse lattice    &                              & $4^4$ \\
near-null vectors & $N_v$                        & $8$ \\
coarse dof / site & $2N_v$ (chirality split)     & $16$ \\
setup smoothing   & $N_{\mathrm{step}}$          & $0$ (unoptimised) \\
block orthog.     &                              & batched CGS, single pass \\
smoother          &                              & SAP, 4 sweeps, tol $10^{-3}$ \\
coarse solver     &                              & GMRES, $\le2$ iter, tol $10^{-2}$ \\
outer solver      &                              & FGMRES, restart $N_M=30$ \\
grid transfer / coarse prec. &                   & single (\texttt{float}) \\
\bottomrule
\end{tabular}
\caption{Multigrid setup and cycle parameters.}
\label{tab:mg}
\end{table}

\section{Reference solvers}
\label{sec:reference}
We compare against two baselines that solve the \emph{same} physical propagator,
with the \emph{same} chiral source/sink construction of Sec.~\ref{sec:projection}:
\begin{itemize}
\item \textbf{Even-odd CGNE} on the normal equation $\Ddw\dg\Ddw$ --- the
production domain-wall solver, run in standard precision.
\item \textbf{Unpreconditioned GMRES on $\Aop$} --- restarted GMRES applied to the
identical target operator \eqref{eq:A} with the identical left-preconditioned
right-hand side, but \emph{without} the multigrid preconditioner. AMG and this
baseline differ \emph{only} by the V-cycle, so their comparison isolates the
preconditioner. Both wrap the inner $\Aop$-solve in the same outer
iterative-refinement on the physical residual (each step solves
$\Aop\,\delta=\PL\,r_{\mathrm{phys}}$, $x\!+\!=\!\delta$,
$r_{\mathrm{phys}}=b-\Ddw x$), so both reach the same physical tolerance.
\end{itemize}

\section{Numerical results}
\label{sec:results}

\paragraph{Setup.} Parameters of Tables~\ref{tab:params} and \ref{tab:mg}; single
NVIDIA RTX~3080. The two operator-$\Aop$ solvers (AMG, GMRES-on-$\Aop$) use the
\emph{identical} left-preconditioned right-hand side. All three solvers are driven
to a comparable \emph{physical} residual $\norm{b-\Ddw x}/\norm{b}\sim10^{-12}$,
so the iteration counts in Table~\ref{tab:results} are residual-matched.

\begin{table}[h]
\centering
\begin{tabular}{lccc}
\toprule
Solver & operator & outer iterations & $\norm{b-\Ddw x}/\norm{b}$ \\
\midrule
\textbf{AMG} (FGMRES + V-cycle) & $\Aop$ & \textbf{210} & $1.3\times10^{-12}$ \\
GMRES on $\Aop$ (no preconditioner) & $\Aop$ & $1202$ & $7.0\times10^{-13}$ \\
Even-odd CGNE (production) & $\Ddw\dg\Ddw$ & $\approx 333$ & $2.4\times10^{-12}$ \\
\bottomrule
\end{tabular}
\caption{Residual-matched iteration counts, $12^4$ $\beta=5.7$, $m_q=0.005$, all
at physical residual $\sim10^{-12}$. Against the identical-operator,
identical-right-hand-side GMRES baseline the V-cycle reduces the outer iteration
count by $\approx5.7\times$; against the production even-odd CGNE by
$\approx1.6\times$. The pseudoscalar correlators agree: the $\Aop$-based and CGNE
$C_{PP}(t)$ coincide to solver tolerance.}
\label{tab:results}
\end{table}

\paragraph{Reading the numbers honestly.}
Three caveats frame Table~\ref{tab:results}.
(i) The clean comparison is AMG versus unpreconditioned GMRES on $\Aop$: same
operator, same right-hand side, same physical tolerance, differing only by the
V-cycle ($5.7\times$). The CGNE column is the practical reference on a different
operator ($\Ddw\dg\Ddw$).
(ii) An outer AMG iteration (one V-cycle: SAP sweeps $+$ coarse solve) is much
more expensive than a CGNE iteration, so the $1.6\times$ iteration advantage over
CGNE does \emph{not} translate into a wall-clock advantage on this small, coarse
problem --- and is not claimed to.
(iii) Most importantly, the near-null setup is unoptimised ($N_{\mathrm{step}}=0$,
$N_v=8$, two grid levels, single consumer GPU). A tuned setup is expected to
lower the AMG count substantially.

\paragraph{Where the value is.}
The textbook value of multigrid is that the outer iteration count is essentially
independent of the quark mass and the volume, whereas CGNE grows as $m_q\to0$ and
with the volume. Quantifying that scaling --- with an optimised near-null setup,
on larger volumes and smaller $m_q$, on hardware with the memory and throughput to
support it --- is the natural next step, and is the case for moving the study to a
supercomputer. The present run establishes correctness (matching $C_{PP}$ and
$m_{\mathrm{res}}$ across all three solvers) and a residual-matched baseline on a
single RTX~3080.

\subsection{Batched, single-precision propagator through \texorpdfstring{$\Aop$}{A}}
\label{sec:batched}

The twelve colour--spin columns of the point-source propagator are solved through
$\Aop$ as a single multiple-right-hand-side (MRHS) block: a block-FGMRES outer
solver preconditioned by the batched V-cycle (batched $\Aop$-apply, restriction
and prolongation cast as small dense GEMMs, and a batched coarse block-FGMRES on
the Galerkin operator), with the physical propagator recovered by an outer
residual refinement built on the left-preconditioner. Because $\Aop=\PL\Ddw$, the
identity $\Aop^{-1}\PL=\Ddw^{-1}$ holds exactly, so a single refinement returns
the CGNE propagator up to the inner-solve tolerance. The entire V-cycle, the
$\Aop$-apply and the coarse solve run in FP32; double precision enters only the
optional outer residual, selected at runtime. On the single RTX~3080 the twelve
columns are processed in MRHS sub-batches of four --- a memory/throughput knob:
the full twelve-wide batch is throughput-optimal (one block solve, amortised
setup, maximal occupancy) but its block-Krylov basis and refinement scratch
exceed the $10$~GB device, whereas data-centre GPUs run the full batch.

The resulting pseudoscalar correlator reproduces the production even-odd CGNE,
\[
  C_{PP}^{\mathrm{AMG}} = 1.56750\times10^{-1}, \qquad
  C_{PP}^{\mathrm{CGNE}} = 1.56754\times10^{-1},
\]
agreeing to $2.4\times10^{-4}$ --- the FP32 floor of the single-precision inner
solve, well below any physical scale. This upgrades the qualitative
``$C_{PP}$ coincide to solver tolerance'' of Table~\ref{tab:results} to a
quantitative, fully-batched, FP32 check.

\paragraph{The \texorpdfstring{$\alpha$}{alpha}-invariance check.}
Carrying the Neff $\alpha$-form through the \emph{physical} operator (and hence,
by the Galerkin construction, through the coarse operator) should leave the
four-dimensional pion invariant while shifting the residual mass as
$1/\alpha^2$. Repeating the batched solve at $\alpha=1$ and $\alpha=0.5$ on the
same solver (Table~\ref{tab:alpha}) confirms both: the pion correlator is
invariant to $1.8\times10^{-4}$ (the FP32 floor), while $m_{\mathrm{res}}$ grows
by $\approx3.6\times$, broadly consistent with the factor $1/\alpha^2=4$ (the
departure is attributable to the FP32 inner-solve floor and the small,
artificially-conditioned test lattice). This validates the $\alpha$-form physics
and confirms that $\alpha$ propagates correctly into the Galerkin coarse operator
--- the prerequisite for using $\alpha$ as the coarse-space spectral knob.

\begin{table}[h]
\centering
\begin{tabular}{lcc}
\toprule
batched AMG & $C_{PP}$ & $m_{\mathrm{res}}$ \\
\midrule
$\alpha=1.0$ & $1.56750\times10^{-1}$ & $9.66\times10^{-2}$ \\
$\alpha=0.5$ & $1.56779\times10^{-1}$ & $3.50\times10^{-1}$ \\
\bottomrule
\end{tabular}
\caption{$\alpha$-invariance on the batched AMG solver, $12^4$, $\beta=5.7$,
$m_q=0.005$. The pion $C_{PP}$ is invariant to the FP32 floor; the residual mass
shifts $\approx3.6\times$, consistent with the $1/\alpha^2$ scaling of the
$\alpha$-form.}
\label{tab:alpha}
\end{table}

\section{Future work}
\label{sec:future}

The present note establishes correctness and a residual-matched baseline. The
following programme of validation is ordered so that each inexpensive step
de-risks the more expensive ones that follow; item~1 is a prerequisite for all
of them.

\begin{enumerate}
\item \textbf{Near-null subspace setup.} The run reported here uses
$N_{\mathrm{step}}=0$ setup smoothing and $N_v=8$ vectors. The first task is to
restore an adaptive setup ($N_{\mathrm{step}}>0$, larger $N_v$) and measure the
reduction in outer iteration count, so that the AMG is assessed at its genuine
optimum rather than the present deliberately minimal configuration.

\item \textbf{Quark-mass scaling.} Measure the outer iteration count (and the
fine-operator application count) of the AMG, of unpreconditioned GMRES on $\Aop$,
and of even-odd CGNE as a function of $m_q\in\{0.01,0.005,0.002,0.001\}$ at fixed,
residual-matched tolerance. The defining property of multigrid is a near
mass-independent iteration count, set against the steep growth of CGNE as
$m_q\to0$; this scaling, not a single-point ratio, is the quantitative case for
the method.

\item \textbf{The $\alpha$ parameter as a spectral knob.} The Neff $\alpha$-form
of the M\"obius operator (arXiv:2602.14454) leaves the four-dimensional pion
invariant while shifting the residual mass as $1/\alpha^2$ and compressing the
spectral range of $\Aop$ (optimal $\alpha\approx0.5$). Boyle and Yamaguchi
(arXiv:2103.05034) identify the wide coarse-space spectral range $[m_l,64]$ of the
PV-preconditioned operator as a limiting factor for the MG-PV class; $\alpha$ is
precisely a knob on that range. We will apply the optimal $\alpha$ inside the
V-cycle coarse and smoother operators only, holding the physical operator at
$\alpha=1$, and measure the reduction in outer iterations (with $C_{PP}$ and
$m_{\mathrm{res}}$ unchanged as the $\alpha$-invariance check). This lever is
absent from all existing domain-wall multigrid studies.

\item \textbf{Multiword / FP32 precision in the multigrid path.} All published
domain-wall multigrid runs use double precision throughout. We will carry the
V-cycle and coarse-grid work in float--float (double-word) arithmetic and verify
that the true residual reaches the FP64 floor at FP32 storage and bandwidth cost.
This precision reach --- recovering double-precision accuracy from
single-precision hardware --- is the principal novel direction, and it interacts
with the $\alpha$ knob through the spectral range the coarse polynomial must
resolve.

\item \textbf{Multiple right-hand sides (MRHS).} Following Boyle
(arXiv:2409.03904), batch the twelve colour--spin columns of the propagator (and,
in the HMC context, multiple Hasenbusch factors) through $\Aop$ and the V-cycle
with a block-FGMRES outer solver and a batched (GEMM) coarse operator. This makes
the coarse-space cost --- the dominant expense of the five-dimensional
PV-preconditioned coarse space, and the factor that rendered the MG-PV class
uncompetitive in earlier comparisons --- sub-dominant, and adds a block-Krylov
algorithmic gain on top of the throughput gain.

\item \textbf{Combined MRHS $\times$ multiword $\times$ FP32.} Batched
single-precision GEMM on tensor hardware, with multiword arithmetic restoring the
precision, applied to the full V-cycle. This combination --- a precise,
PV-preconditioned, positive-operator domain-wall multigrid on batched
single-precision tensor cores --- is unexplored, and is attempted only once
items~4 and 5 succeed in isolation.
\emph{Status update.} The MRHS batched propagator and the FP32 V-cycle (in
single-precision, not yet multiword form) are now implemented and validated; see
Section~\ref{sec:batched}. The $\alpha$ lever has been exercised as a physics
check (Table~\ref{tab:alpha}); applying it to the preconditioner only, for the
iteration-count reduction, remains open.

\item \textbf{Scale-up validation.} Repeat the comparison on larger volumes and
smaller quark masses, on data-centre GPUs (A100/H100) whose memory and
batched-GEMM throughput the method is designed to exploit. The single consumer GPU
used here understates the achievable gain; this step is the quantitative basis for
a production deployment.
\end{enumerate}

\end{document}
